% © 2026 Ing. David Jaroš — CC BY-NC-ND 4.0
%
% This work is licensed under a Creative Commons Attribution-NonCommercial-NoDerivatives
% 4.0 International License (CC BY-NC-ND 4.0).
%
% Δd = 0.405 — Explaining the gap B₀ → B_base
% Track: CORE — α derivation — v59 baseline
% Date: 2026-03-08

\ifdefined\INCLUDEMODE
  % Being included in another document — skip preamble
\else
  \documentclass[11pt,a4paper]{article}
  \usepackage{amsmath,amssymb,amsthm}
  \usepackage{hyperref}

  \newtheorem{theorem}{Theorem}
  \newtheorem{lemma}{Lemma}
  \newtheorem{proposition}{Proposition}
  \newtheorem{conjecture}{Conjecture}
  \theoremstyle{definition}
  \newtheorem{gap}{Gap}
  \theoremstyle{remark}
  \newtheorem{remark}{Remark}

  \title{$\Delta d = 0.405$ — Explaining the Transition $B_0 \to B_{\mathrm{base}}$\\
    \large Approaches C1 (Heat Kernel Curvature) and C2 (Mode Interference)}
  \author{Ing.\ David Jaroš}
  \date{2026-03-08 (v59)}

  \begin{document}
  \maketitle
\fi

\section{Approach C: Explaining $\Delta d = 0.405$ — the gap $B_0 \to B_{\mathrm{base}}$}
\label{sec:delta_d_main}

\textbf{Overall status: [OPEN GAP — both approaches below lead to dead ends].}
No circular use of $\alpha^{-1} = 137$ or $n^* = 137$ anywhere in this document.
$N_{\mathrm{eff}} = 12$ and $\dim_\mathbb{R}(\mathrm{Im}\,\mathbb{H}) = 3$ are fixed inputs (both proved).

%--------------------------------------------------------------------
\subsection{Section 1: Problem Formulation}
\label{sec:formulation}
%--------------------------------------------------------------------

From the v59 baseline the following quantities are established:

\begin{center}
\begin{tabular}{lll}
  \hline
  Quantity & Value & Status \\
  \hline
  $B_0 = 2\pi N_{\mathrm{eff}}/3$ & $25.133$ & \textbf{[PROVED]} (one-loop, flat space) \\
  $B_{\mathrm{base}} = N_{\mathrm{eff}}^{3/2}$ & $41.569$ & [CONJECTURED] \\
  $d_{\mathrm{eff}}(B_0)$ & $2.595$ & \textbf{[PROVED]} (from one-loop $B_0$) \\
  $d_{\mathrm{eff}}(B_{\mathrm{base}})$ & $3.000$ & algebraic identity \\
  $\Delta d = d - d_{\mathrm{eff}}(B_0)$ & $0.405$ & \textbf{[OPEN]} \\
  \hline
\end{tabular}
\end{center}

where the effective dimension is defined by
$d_{\mathrm{eff}}(B) := 2\log B / \log N_{\mathrm{eff}}$,
so that $B = N_{\mathrm{eff}}^{d_{\mathrm{eff}}/2}$
(see Remark~\ref{rem:d_eff} in \texttt{b\_base\_hausdorff.tex}).

The algebraic ratio is an exact identity:
\begin{equation}
  \frac{B_{\mathrm{base}}}{B_0}
    = \frac{N_{\mathrm{eff}}^{3/2}}{2\pi N_{\mathrm{eff}}/3}
    = \frac{3}{2\pi}\,\sqrt{N_{\mathrm{eff}}}
    = \frac{3}{2\pi}\,\sqrt{12}
    \approx 1.654.
  \label{eq:ratio_identity}
\end{equation}

\textbf{Question:} What physical mechanism lifts $d_{\mathrm{eff}}$ from $2.595$
(proved one-loop flat-space value) to $3.000$
(algebraic value $\dim_\mathbb{R}(\mathrm{Im}\,\mathbb{H})$)?
Two candidate mechanisms are investigated below.

%--------------------------------------------------------------------
\subsection{Section 2: Approach C1 — Heat Kernel Curvature Correction on $\mathrm{Im}(\mathbb{H})$}
\label{sec:C1_curvature}
%--------------------------------------------------------------------

\textbf{Hypothesis.}
The proved $B_0$ is computed in flat space $\mathbb{R}^d$.
However, $\mathrm{Im}(\mathbb{H})$ carries the structure of the Lie algebra
$\mathfrak{su}(2)$; exponentiating gives $\mathrm{SU}(2) \cong S^3$, which has
non-zero scalar curvature.
Could the Seeley--DeWitt curvature correction lift $B_0 \to B_{\mathrm{base}}$?

\subsubsection{Step 1: Heat kernel representation of $B_0$}

The one-loop contribution to the renormalisation coefficient $B$ can be written as
\begin{equation}
  B_0 \;=\; \int_0^{1/\Lambda^2} \frac{dt}{t}\;
           \mathrm{Tr}\,K_{\mathrm{flat}}(t),
  \qquad
  K_{\mathrm{flat}}(t,x,x) = (4\pi t)^{-d/2},
  \label{eq:B0_heatkernel}
\end{equation}
where $\Lambda$ is the UV cutoff and $d = 3$.
This integral is UV divergent; the regularised value $B_0 = 2\pi N_{\mathrm{eff}}/3$
is obtained after including the gauge-structure factor $N_{\mathrm{eff}}$ and the
$2\pi/d$ prefactor from the one-loop calculation
(\texttt{b\_base\_spinor\_approach.tex}, \texttt{b\_base\_hausdorff.tex}).

\subsubsection{Step 2: Seeley--DeWitt correction from curvature of $\mathrm{Im}(\mathbb{H})$}

The standard heat kernel expansion on a $d$-dimensional Riemannian manifold $M$
with scalar curvature $R$ is
\begin{equation}
  K(t,x,x)
    = (4\pi t)^{-d/2}\bigl[1 + \tfrac{R}{6}\,t + O(t^2)\bigr].
  \label{eq:seeley_dewitt}
\end{equation}
The curvature correction shifts $B$ by
\begin{equation}
  \Delta B_{\mathrm{curv}}
    = \frac{R_{\mathrm{Im}\mathbb{H}}}{6}
      \int_0^{1/\Lambda^2} \frac{dt}{t}\;(4\pi t)^{-d/2}\cdot t
    = \frac{R_{\mathrm{Im}\mathbb{H}}}{6}\;
      \int_0^{1/\Lambda^2} dt\;(4\pi)^{-3/2}\,t^{-3/2}.
  \label{eq:DeltaB_curv}
\end{equation}

For $d = 3$, the integrals evaluate to:
\begin{align}
  B_0 &\;\propto\; \int_0^{1/\Lambda^2} dt\;t^{-5/2}
       \;=\; \tfrac{2}{3}\,\Lambda^3
       \quad (\text{leading UV divergence}),
       \label{eq:B0_UV} \\[4pt]
  \Delta B_{\mathrm{curv}} &\;\propto\;
       \frac{R_{\mathrm{Im}\mathbb{H}}}{6}
       \int_0^{1/\Lambda^2} dt\;t^{-3/2}
    \;=\; \frac{R_{\mathrm{Im}\mathbb{H}}}{6}\cdot 2\Lambda.
       \label{eq:DeltaB_UV}
\end{align}
Therefore:
\begin{equation}
  \frac{\Delta B_{\mathrm{curv}}}{B_0}
    \;\propto\; \frac{R_{\mathrm{Im}\mathbb{H}}}{6}\cdot \frac{2\Lambda}{(2/3)\Lambda^3}
    = \frac{R_{\mathrm{Im}\mathbb{H}}}{2\Lambda^2}
    \;\xrightarrow{\Lambda \to \infty}\; 0.
  \label{eq:ratio_curv}
\end{equation}

\textbf{Key result:} the Seeley--DeWitt curvature correction is \emph{subleading
in the UV limit} — it vanishes as $\Lambda \to \infty$.
The ratio $\Delta B_{\mathrm{curv}} / B_0 \propto \Lambda^{-2} \to 0$,
regardless of the value of $R_{\mathrm{Im}\mathbb{H}}$.
This is a general feature of Seeley--DeWitt corrections in $d \ge 2$:
the curvature term enters at one order lower in $\Lambda$ than the leading divergence.

\subsubsection{Step 3: Required curvature (numerical check, massive regularisation)}

If one instead uses a massive regulator $M$ (replacing $\Lambda$) and evaluates
$B$ at $t = 1/M^2$:
\begin{equation}
  \frac{B_{\mathrm{curved}}}{B_{\mathrm{flat}}}
    = 1 + \frac{R_{\mathrm{Im}\mathbb{H}}}{6\,M^2}
    \stackrel{!}{=} \frac{B_{\mathrm{base}}}{B_0} \approx 1.654.
  \label{eq:R_required}
\end{equation}
This requires:
\begin{equation}
  R_{\mathrm{Im}\mathbb{H}} = 0.654 \times 6 M^2 = 3.924\,M^2.
  \label{eq:R_numerical}
\end{equation}
For the natural normalization $M = 1$ (unit quaternion scale $|i| = |j| = |k| = 1$):
$R_{\mathrm{Im}\mathbb{H}} = 3.924$.

\noindent\textbf{Comparison with $\mathfrak{su}(2)$ geometry:}
$\mathrm{Im}(\mathbb{H})$ as a vector space is flat $\mathbb{R}^3$ (curvature zero).
The curvature arises only when we identify $\exp(\mathrm{Im}(\mathbb{H})) = \mathrm{SU}(2) \cong S^3$.
For the bi-invariant metric on $\mathrm{SU}(2)$ induced by the quaternion inner product
$\langle e_a, e_b \rangle = \delta_{ab}$:
\begin{equation}
  R_{S^3}(r=1) = \frac{6}{r^2}\Big|_{r=1} = 6.
  \label{eq:R_S3}
\end{equation}
The value $R = 6$ (for the unit sphere) does not match the required $R = 3.924$
from~\eqref{eq:R_numerical}.
To match, one would need $r = \sqrt{6/3.924} \approx 1.237$,
which has no natural algebraic meaning.

\begin{gap}[Approach C1 — Curvature correction]
\label{gap:C1}
  \textbf{[DEAD END].}
  Two independent reasons:
  \begin{enumerate}
    \item[(i)] For UV-divergent integrals (power-law cutoff $\Lambda$), the
      Seeley--DeWitt correction is subleading: $\Delta B / B_0 \propto \Lambda^{-2} \to 0$.
      It cannot produce a finite correction factor $\approx 0.654$.
    \item[(ii)] For a massive regulator $M = 1$ (natural quaternion unit),
      the required curvature is $R = 3.924$,
      while the natural $S^3$ value is $R = 6$ (unit radius),
      or zero (flat $\mathbb{R}^3$).
      No normalization of $\mathfrak{su}(2)$ gives $R = 3.924$ without a free parameter.
  \end{enumerate}
  The hypothesis that $\Delta d = 0.405$ arises from the Seeley--DeWitt curvature
  correction on $\mathrm{Im}(\mathbb{H})$ is \textbf{falsified by the power-counting argument}.
\end{gap}

%--------------------------------------------------------------------
\subsection{Section 3: Approach C2 — Mode Pair Interference}
\label{sec:C2_interference}
%--------------------------------------------------------------------

\textbf{Hypothesis.}
The factor $B_{\mathrm{base}}/B_0 = (3/2\pi)\sqrt{N_{\mathrm{eff}}}$
might arise from coherent interference of $\sim\!\sqrt{N_{\mathrm{eff}}}$ mode pairs
out of the $N_{\mathrm{eff}}$ total modes.

\subsubsection{Naive $\sqrt{N_{\mathrm{eff}}}$ coherence}

If $\sqrt{N_{\mathrm{eff}}}$ modes contribute coherently, one might expect
$B_{\mathrm{eff}} = B_0 \cdot \sqrt{N_{\mathrm{eff}}}$.
Numerically:
\begin{equation}
  B_0 \cdot \sqrt{N_{\mathrm{eff}}} = 25.133 \times \sqrt{12} = 25.133 \times 3.464 \approx 87.05
  \;\ne\; B_{\mathrm{base}} = 41.569.
  \label{eq:naive_sqrt_fails}
\end{equation}
\textbf{This version fails immediately.}
The naive $\sqrt{N_{\mathrm{eff}}}$ scaling overshoots by a factor of $2.09$.

\subsubsection{Decomposing the ratio}

The algebraic identity \eqref{eq:ratio_identity} decomposes as:
\begin{equation}
  \frac{B_{\mathrm{base}}}{B_0}
  = \frac{3}{2\pi}\,\sqrt{N_{\mathrm{eff}}}
  = \underbrace{\frac{d}{2\pi}}_{\text{from }B_0\text{ formula}}
    \cdot \underbrace{\sqrt{N_{\mathrm{eff}}}}_{\text{factor to explain}}
  = \frac{N_{\mathrm{eff}}^{3/2}}{(2\pi/d)\,N_{\mathrm{eff}}}.
  \label{eq:ratio_decomposed}
\end{equation}
In other words, $B_{\mathrm{base}}/B_0 = (3/2\pi)\sqrt{N_{\mathrm{eff}}}$ is precisely
the statement $B_{\mathrm{base}} = N_{\mathrm{eff}}^{3/2}$ rewritten using $B_0 = (2\pi/d)\,N_{\mathrm{eff}}$.
Any ``physical interpretation'' of the $\sqrt{N_{\mathrm{eff}}}$ factor would therefore
need to explain $B_{\mathrm{base}} = N_{\mathrm{eff}}^{3/2}$ from first principles,
which is precisely Gap~(a) in \texttt{b\_base\_hausdorff.tex}.

\subsubsection{Can $\sqrt{N_{\mathrm{eff}}}$ arise from pair counting?}

The number of distinct unordered pairs of $N_{\mathrm{eff}}$ modes is
$\binom{N_{\mathrm{eff}}}{2} = N_{\mathrm{eff}}(N_{\mathrm{eff}}-1)/2 = 66$ (for $N_{\mathrm{eff}} = 12$).
If only a fraction $\sim 1/\sqrt{N_{\mathrm{eff}}}$ of pairs contributes coherently:
$66 / \sqrt{12} \approx 19$, which has no obvious connection to the prefactor $3/2\pi$.

Alternatively, if the coherent sum scales as $N_{\mathrm{eff}}^{1/2}$ (random-walk-like),
then $B = B_{\mathrm{single}} \cdot N_{\mathrm{eff}}^{1/2}$ requires
$B_{\mathrm{single}} = B_{\mathrm{base}} / N_{\mathrm{eff}}^{1/2} = 12^{3/2}/12^{1/2} = 12$.
But $B_0/N_{\mathrm{eff}} = 25.133/12 \approx 2.09 \ne 12$,
so there is no consistent single-mode contribution.

\begin{gap}[Approach C2 — Mode pair interference]
\label{gap:C2}
  \textbf{[DEAD END].}
  \begin{enumerate}
    \item[(i)] Naive $\sqrt{N_{\mathrm{eff}}}$ coherence:
      $B_0 \cdot \sqrt{N_{\mathrm{eff}}} \approx 87.05 \ne 41.57$.
      Off by factor $2.09$.
    \item[(ii)] The ratio $B_{\mathrm{base}}/B_0 = (3/2\pi)\sqrt{N_{\mathrm{eff}}}$
      is an algebraic identity following directly from $B_{\mathrm{base}} = N_{\mathrm{eff}}^{3/2}$
      and $B_0 = (2\pi/d)N_{\mathrm{eff}}$.
      Explaining it via mode interference is equivalent to explaining the conjecture
      $B_{\mathrm{base}} = N_{\mathrm{eff}}^{3/2}$ itself — it does not provide an independent derivation.
    \item[(iii)] No pairing or coherence prescription yields the factor $(3/2\pi)$
      without a free parameter.
  \end{enumerate}
  The mode-pair-interference hypothesis does not provide a new derivation
  and reduces to restating the original conjecture.
\end{gap}

%--------------------------------------------------------------------
\subsection{Section 4: Numerical Control — for which $N_{\mathrm{eff}}$ is $d_{\mathrm{eff}}(B_0) = 3$?}
\label{sec:numerical_control}
%--------------------------------------------------------------------

The condition $d_{\mathrm{eff}}(B_0) = d = 3$ is:
\begin{equation}
  \frac{2\log(2\pi N/3)}{\log N} = 3
  \;\;\Longrightarrow\;\;
  \bigl(2\pi N/3\bigr)^2 = N^3
  \;\;\Longrightarrow\;\;
  \frac{4\pi^2}{9} = N.
  \label{eq:N_star_condition}
\end{equation}
Hence:
\begin{equation}
  N^* \;:=\; \frac{4\pi^2}{9} \;\approx\; 4.386.
  \label{eq:N_star_value}
\end{equation}
Verification: $d_{\mathrm{eff}}(B_0)\big|_{N=4.386}
= 2\log(2\pi\cdot 4.386/3)/\log(4.386)
= 2\log(9.189)/\log(4.386)
\approx 2(2.218)/1.477 \approx 3.003$. \checkmark

\medskip
\noindent
Comparison for selected values of $N_{\mathrm{eff}}$:

\begin{center}
\begin{tabular}{cccc}
  \hline
  $N_{\mathrm{eff}}$ & $B_0 = 2\pi N/3$ & $d_{\mathrm{eff}}(B_0)$ & $\Delta d = 3 - d_{\mathrm{eff}}(B_0)$ \\
  \hline
  $4.386 = N^*$ & $9.188$ & $3.000$ & $0.000$ \\
  $4$ & $8.378$ & $3.067$ & $-0.067$ \\
  $6$ & $12.566$ & $2.825$ & $+0.175$ \\
  $8$ & $16.755$ & $2.711$ & $+0.289$ \\
  $12$ & $25.133$ & $2.595$ & $+0.405$ \quad {\small $\leftarrow$ UBT} \\
  $24$ & $50.265$ & $2.465$ & $+0.535$ \\
  \hline
\end{tabular}
\end{center}

\begin{remark}[$\Delta d = 0.405$ is specific to $N_{\mathrm{eff}} = 12$]
\label{rem:delta_d_specific}
  The gap $\Delta d = d - d_{\mathrm{eff}}(B_0)$ depends on $N_{\mathrm{eff}}$:
  it equals zero only at $N^* = 4\pi^2/9 \approx 4.386$,
  and grows monotonically with $N_{\mathrm{eff}}$ for $N_{\mathrm{eff}} > N^*$.
  For the UBT value $N_{\mathrm{eff}} = 12$, $\Delta d = 0.405$.
  This number is not a free parameter: it is a \emph{derived consequence}
  of $N_{\mathrm{eff}} = 12$ (proved from $\mathbb{C}\otimes\mathbb{H}$ algebra)
  and $d = 3$ (proved from $\dim_\mathbb{R}(\mathrm{Im}\,\mathbb{H}) = 3$).
  The physical mechanism that enforces $d_{\mathrm{eff}} \to d$ despite $N_{\mathrm{eff}} \ne N^*$
  remains the core open question.
\end{remark}

%--------------------------------------------------------------------
\subsection{Section 5: Conclusion — Status Summary}
\label{sec:conclusion}
%--------------------------------------------------------------------

\begin{center}
\begin{tabular}{lll}
  \hline
  Approach & Result & Status \\
  \hline
  C1: Seeley--DeWitt curvature on $\mathrm{Im}(\mathbb{H})$
    & Correction $\propto \Lambda^{-2} \to 0$
    & \textbf{[DEAD END]} \\
  C2: Mode pair interference ($\sqrt{N_{\mathrm{eff}}}$ coherence)
    & Algebraic identity restatement
    & \textbf{[DEAD END]} \\
  \hline
\end{tabular}
\end{center}

Gap~(a) from \texttt{b\_base\_hausdorff.tex} therefore remains \textbf{[OPEN]}.

\medskip\noindent\textbf{What was learned:}
\begin{itemize}
  \item $\Delta d = 0.405$ is a derived numerical consequence of $N_{\mathrm{eff}} = 12$
    and $d = 3$ (both proved); it is not a free parameter
    (Remark~\ref{rem:delta_d_specific}).
  \item A standard Seeley--DeWitt correction on the mode space cannot lift
    $d_{\mathrm{eff}}$ by a finite amount independent of the UV cutoff,
    because curvature corrections are subleading in the UV.
  \item The algebraic ratio $(3/2\pi)\sqrt{N_{\mathrm{eff}}}$ is an identity,
    not an independent mechanism.
  \item The open question is: what renormalisation effect or geometric structure
    constrains $d_{\mathrm{eff}}(B)$ to be exactly the algebraic dimension $d = 3$?
    Possible directions not yet explored: zeta-function regularisation of
    $\mathrm{Spec}(-\nabla^2|_{\mathrm{Im}(\mathbb{H})})$;
    index theorem constraints; non-perturbative effects.
\end{itemize}

%--------------------------------------------------------------------
\subsection{Section 6: In Plain Language}
\label{sec:plain}
%--------------------------------------------------------------------

\textbf{In plain language:}
The proved one-loop calculation gives $B_0 \approx 25.13$, which corresponds to
an ``effective dimension'' of $2.60$ rather than the algebraic value $3.00$.
The gap $\Delta d = 0.405$ is the amount by which some mechanism must lift this
effective dimension to recover the conjectured $B_{\mathrm{base}} = N_{\mathrm{eff}}^{3/2} \approx 41.57$.
Two candidate mechanisms were investigated: (C1) a curvature correction from the
quaternionic geometry of the mode space, and (C2) coherent interference of mode pairs.
Both turned out to be dead ends — C1 because curvature corrections in quantum field
theory are subleading (negligible) in the ultraviolet limit, and C2 because the
$\sqrt{N_{\mathrm{eff}}}$ factor is merely an algebraic restatement of the conjecture itself.
The gap remains open; explaining $\Delta d = 0.405$ from first principles is
equivalent to completing the proof of $B_{\mathrm{base}} = N_{\mathrm{eff}}^{3/2}$.

\ifdefined\INCLUDEMODE\else
\end{document}
\fi
